\subsection{What the apparatus does}
\label{sec:analysis}

Both the standings of \S\ref{sec:results} and the behaviour of \S\ref{sec:behavior} are read
against a fixed apparatus, and several properties of that apparatus move the
reported number by more than adjacent mid-table systems differ from one another. Four of those properties are worth making explicit: how the
checkpoint rule changes the headline without touching an agent, how many of the resulting wins the
evaluator can actually resolve, what a unit of exploration costs, and the two tasks on which the
repository's own metric stops measuring what it was chosen to stand for.
A fourth --- what fails and why --- belongs with the population it removes and was given in
\S\ref{sec:lifecycle}.

\subsubsection{The checkpoint rule decides which artifact is reported}
\label{sec:ckpt}

Each configuration keeps up to three checkpoints and is scored on the best of them under the task's
direction. It changes the reported artifact for most of the suite and the
per-artifact scores very little, yet it moves the headline. Best
differs from final in $145$ of $271$ scored configurations ($53.5\%$), and the winning checkpoint is
the first about as often as the last, so this is not a late-training effect. The mechanism is
threshold-straddling: most configurations land within a fraction of a percent of the recipe, so a
per-artifact change too small to matter on its own flips the binary win or loss for many of them. Its consequence for the
aggregate is large: reading the nine scalar-metric tasks on final rather than best-of-three
checkpoints moves $19$ of $243$ configurations across the line, a swing of $7.8$ points --- more than separates
adjacent mid-table systems. A benchmark reporting best-of-$k$ and one reporting final checkpoints would agree
about the typical configuration on these runs and disagree about the headline, which is why we report
the rule alongside the number. One further property of the rule is worth stating because it is not
symmetric: the harness fixes \emph{how many} checkpoints are reduced, but not \emph{what occupies}
them. The gallery records agents filling the three slots with a portfolio of distinct recipes rather
than three points on one training curve --- a three-member sweep over the preference temperature,
three pruning calibration seeds, the best three of eleven restarts chosen by a candidate-owned
validator. Selection over a curve and selection over a portfolio are different things, and the rule
permits both.

\subsubsection{How many wins survive the evaluator's own noise}
\label{sec:noise}

The same threshold-straddling has a second consequence the win rate hides. Seven of the ten metrics are
means over a countable evaluation set --- images, prompts, problems, boards, preference pairs --- and
the frozen evaluator records a standard error for each artifact alongside its score. On the other three
it does not, and cannot: a perplexity and a test NLL are deterministic functions of a model and a fixed
corpus rather than sample means. So of the $152$ configurations that beat the shipped recipe on the
nine scalar-metric tasks, $101$ carry a testable margin and $51$ (graph generation and one-shot
pruning) do not.

Of those $101$, \textbf{$42\%$ beat the recipe by more than twice the evaluator's
standard error} and $74\%$ by more than one. We do not hold a standard error for the reference itself;
inflating the artifact's own by a factor of $1.41$ as a stand-in for the error of the difference, a
deliberately strict substitute, leaves $38\%$ above two. The rest are wins the count records and the measurement
cannot separate from noise. The concentration matters more than the pooled figure: on diffusion
fine-tuning $22$ of $26$ testable wins clear two standard errors and on agentic RL $15$ of $17$ do,
while on weight averaging only $1$ of $28$ does and on on-policy distillation \emph{none} of $14$ does.
Weight averaging is one of the two tasks whose recipe \S\ref{sec:results} already
calls a soft reference, and this is the same softness measured a second way. On-policy distillation is
not on that list: the win rate reports it at $14/25$, yet its result rests entirely on differences the
evaluator cannot resolve. Its own task package is the only one
of the ten to publish an uncertainty on the reference, $\pm0.0516$ against a recipe of $0.4271$, which
is the same fact seen from the other side.

We report this rather than filter on it. A margin inside the standard error is not a demonstration that
the recipe was not improved, and the $51$ untestable wins are untestable, not failed; applying a
significance filter to the headline would trade one arbitrary convention for another. But a benchmark
whose pooled rate is carried by tasks where most margins sit inside the evaluator's own resolution is
reporting something weaker than the rate suggests, and the per-artifact errors are released so that a
reader can apply whatever threshold they prefer.

\subsubsection{What a unit of exploration costs}
\label{sec:cost}

\begin{figure}[t]
\centering
\includegraphics[width=\linewidth]{fig_cost}
\caption{\textbf{Cost buys almost no score --- across systems as well as within one.} Each point is one
\emph{(model, effort)} cell --- median exploration cost per configuration (log $x$, colour = system,
size = effort) against win rate over the shipped recipe ($y$). The complement to
Figure~\ref{fig:effort}, which sweeps effort inside one harness: here the whole field is on one axis.
The strongest system, Claude Opus~5, is neither the cheapest nor the most expensive: median cost per
configuration ranges roughly tenfold across the six systems while capability barely tracks it. The ordinary-least-squares fit (with
its $95\%$ band) has a slope of only $+4.5$ win-rate points per tenfold increase in cost, and the
correlation explains about $5\%$ of the variance across the $28$ cells that carry a cost (Pearson
$r=0.21$): the cheapest third of cells (median \$3/cfg, $55\%$ win) reaches
nearly the same score as the priciest (\$35/cfg, $60\%$). Cost is read per configuration so grid depth
does not inflate it; Kimi~K3 is a single max-effort cell.}
\label{fig:cost}
\end{figure}

Aggregated over configurations --- grids of unequal depth, so not a per-system-comparable total --- the
whole evaluation cost about \$5{,}200 in exploration API calls at each provider's official rates, kept
distinct from formal GPU training and the hidden evaluator's compute (Table~\ref{tab:systems}). Two facts stand out (Figure~\ref{fig:cost}). Within the Codex
harness, cost per configuration climbs steeply with effort while the win rate does not
(\S\ref{sec:effort}), so the marginal dollar buys tokens and probes, not wins. Across systems, the
median cost of one configuration ranges roughly ten-fold --- \$45 for GPT-5.6 Sol down to \$4 for
GPT-5.6 Luna --- while medal scores and win rates vary far less, and the strongest system (Claude
Opus~5, \$18 per configuration) is neither the cheapest nor the most expensive. We compare per
configuration rather than per system because the grids differ in depth: the totals span a wider range
still (\$2{,}692 for Sol against \$59 for Kimi~K3's single-effort grid), but that range is mostly a
statement about how many configurations each grid contains. That cross-system comparison is descriptive rather than causal,
since the systems differ in harness as well as in spend; but the weak coupling between cost and
capability is itself the point, and makes cost a first-class axis for an AI-improvement benchmark.

\subsubsection{The reference you divide by, and the metrics that stop measuring}
\label{sec:validity}

Two further caveats bound every number above. First, ``beating the recipe'' depends on which
measurement of the recipe one uses: the shipped recipe run at its committed step count, or that same
configuration replayed under the agent's own twelve-hour budget. The two disagree by task and neither
is uniformly harder. The replay does two things at once. It corrects the reduction asymmetry, and
overcorrects it: the committed run is too short to yield three checkpoints, so $\rshort$ meets the
agent's best of three with a single point, while the replay is reduced to its own best over the twelve
to twenty-nine checkpoints it retains --- more draws than the agent gets. Neither reference is matched
to the agent exactly; between them the agent is bracketed, and the rate against the replay is the
conservative one. And it is the closest thing the suite has to an
\emph{agent-free control}: it opens exactly the two things the patches most often change --- the
training horizon and the checkpoint cadence --- and changes nothing else, so if the wins were bought by
budget and selection rather than by the agent's edits, they should not survive it. They largely do,
and they do so on a pool that excludes both tasks whose committed recipe is a deliberately naive one.
Read on best checkpoints on both sides, agents clear the budget-matched replay in $87$ of $164$
configurations ($53.0\%$) against $80$ of the same $164$ for the recipe as
committed --- seven more, not fewer --- and that near-tie hides
two opposite movements. On three tasks the longer run is the better artifact and therefore the harder
reference, and the win rate falls --- $93\%\!\to\!61\%$ on diffusion fine-tuning, $96\%\!\to\!44\%$ on
graph generation, $56\%\!\to\!20\%$ on on-policy distillation. On two the recipe degrades once its
horizon is opened, so the reference becomes easier and two tasks agents lose on reverse to wins
($24\%\!\to\!76\%$ and $0\%\!\to\!82\%$); reward modelling does not move. A comparison against a single recipe reference is therefore a comparison against
one measurement choice rather than against a fixed bar, which is why we report the committed recipe and its budget-matched replay
separately rather than collapsing them into one ``beat the recipe'' number
(Appendix~\ref{app:baselines}).

Second, on two of the ten tasks the metric each repository chose for itself stops measuring the
intended objective at the top of its range, which is where a best-of-$k$ benchmark reports from.
Neither metric is ours: freezing real repositories means inheriting the objective
each one already optimizes, and an objective that behaves well in its own paper's operating range need
not behave well once an agent is paid to push it. On diffusion fine-tuning the reward is an aesthetic
predictor whose familiar $1$--$10$ range is the span of its training labels rather than a bound, and
the top artifact reaches $17.67$ against a recipe at $5.53$ while its prompt alignment falls $53\%$ and
its output diversity $87\%$. On agentic RL the top scores come from imitating an optimal solver the
repository ships in its own utilities, so the metric separates recipe classes rather than degrees of
training. Appendix~\ref{app:validity} gives both in full. Neither is an evaluator malfunction and
neither breaks a task contract; both stay in the suite with the qualification attached, and together
they are the strongest evidence that validity instrumentation for this kind of benchmark should be
per-task and blocking rather than advisory.
